% =============================================================================
%  V. METHODS: MODELS AND VALIDATION ARCHITECTURE   
% =============================================================================
\begin{frame}{V. Methods: Models and Validation Architecture}
\begin{block}{Primary model: dynamic discrete-time competing-risks hazard}
  \footnotesize Multinomial pooled logistic regression over person-hours
  (sepsis / discharge / death), natural-spline baseline in elapsed
  hours~\cite{dagostino1990pooled,heyard2019dynamic}. The unpenalised fit is
  completely separated, so a small ridge penalty is required for
  identification.
\end{block}

\vspace{2pt}
{\footnotesize\textbf{Comparators:} gradient-boosted trees (XGBoost), Cox
proportional hazards, NEWS2, qSOFA and SIRS~\cite{evans2021ssc}. XGBoost uses
fixed hyperparameters; the number of boosting rounds is chosen by early stopping
on a 15\% stay-level holdout from the \emph{training} data.\par}

\vspace{4pt}
\begin{center}\renewcommand{\arraystretch}{1.12}\scriptsize
\begin{tabular}{p{2.7cm}p{6.6cm}p{3.6cm}}
\toprule
\textbf{Split} & \textbf{How} & \textbf{Why} \\
\midrule
Internal (primary)   & Time-based: train 2008--2016, test 2017--2019~\cite{guo2022temporal} & Primary internal evaluation \\
Internal (secondary) & Random 75/25 patient-level split (seed 42)                    & Split optimism (H4) \\
Treatment-anchored   & Hours before first culture or antibiotic~\cite{kamran2024evaluation} & Treatment leakage (H3) \\
External             & eICU-CRD v2.0, no retraining~\cite{pollard2018eicu}           & Generalisability (H5) \\
\bottomrule
\end{tabular}
\end{center}
{\scriptsize All splits are patient-level: no patient appears in both
training and test sets.\par}
\end{frame}


% =============================================================================
%  V. METHODS: METRICS AND INFERENCE        
% =============================================================================
\begin{frame}{V. Methods: Metrics and Statistical Inference}
\textbf{Primary metric: adapted PhysioNet-style Utility
Score}~\cite{reyna2020utility}
\[
  \widetilde{U} = \frac{U - U_{\text{inact}}}{U_{\text{opt}} - U_{\text{inact}}}
  \qquad
  \widetilde{U}=1 \text{ perfect},\quad
  \widetilde{U}=0 \text{ no-alert strategy},\quad
  \widetilde{U}<0 \text{ worse than no alerts}
\]
\vspace{-6pt}
\begin{itemize}\itemsep2pt\small
  \item Rewards early true alerts in $[t_{\text{onset}}-6,\,t_{\text{onset}}]$,
        penalises late and false ones.
  \item Reported at the 0.5 cut-off \textbf{and} at its maximum over a
        60-point alert-rate grid. Alert burden is false alerts per true alert,
        compared with the published benchmark of 1.4~\cite{moor2023review}.
  \item AUROC, AUPRC and calibration are reported alongside.
\end{itemize}

\vspace{2pt}
\begin{exampleblock}{Inference and reporting}
\begin{itemize}\itemsep1pt\footnotesize
  \item Stay-level cluster bootstrap, $B = 2{,}000$; H1--H6 Holm-corrected
        ($m = 6$, FWER 0.05); exploratory families Benjamini--Hochberg.
  \item Pre-registered hypothesis set; every post-hoc protocol amendment is
        identified explicitly; TRIPOD+AI and PROBAST+AI
        reporting~\cite{collins2024tripodai,moons2025probastai}.
\end{itemize}
\end{exampleblock}
\end{frame}
